% !TEX root = chapter-standalone.tex

\section{From preorders to posets}
\label{sec:preorders-to-posets}

In a \SY{preorder}, two different elements can be ``as good as each other'': it may happen that $\posela \posAleq \poselb$ and $\poselb \posAleq \posela$, although $\posela \neq \poselb$.
This typically happens when the order compares objects by some attribute, such as price or weight, and two different objects happen to have the same value of that attribute.
Antisymmetry is precisely the requirement that this never happens.
In this section we explain how every preorder can be turned into a poset, by treating elements that are as good as each other as one and the same element.

Let $\posA$ be a \SY{preorder}.
We define a relation $\simeq$ on $\posAset$ by
\begin{equation}
    \label{eq:preorder-equivalence}
    \posela \simeq \poselb
    \quad \text{if and only if} \quad
    (\posela \posAleq \poselb) \booland (\poselb \posAleq \posela),
\end{equation}
and we say that $\posela$ and $\poselb$ are \emph{equivalent} in $\posA$.
In the graph representation of $\posA$, two distinct elements are equivalent exactly when they are connected by a double-headed arrow, as in \cref{fig:pre_order_graph}.

\begin{lemma}
    \label{lem:preorder-equivalence}
    Let $\posA$ be a \SY{preorder}.
    The relation $\simeq$ is an \SY{equivalence relation} on $\posAset$.
    Furthermore, $\posA$ is a \SY{poset} if and only if $\posela \simeq \poselb$ implies $\posela = \poselb$.
\end{lemma}
\begin{proof}
    We check the three conditions of \cref{def:equivalence-relation}.
    \begin{enumerate}
        \item Reflexivity: $\posela \posAleq \posela$ holds by reflexivity of $\posAleq$, and therefore $\posela \simeq \posela$.
        \item Symmetry: the definition \eqref{eq:preorder-equivalence} is symmetric in $\posela$ and $\poselb$.
        \item Transitivity: suppose $\posela \simeq \poselb$ and $\poselb \simeq \poselc$.
        From $\posela \posAleq \poselb$ and $\poselb \posAleq \poselc$ we get $\posela \posAleq \poselc$, and from $\poselc \posAleq \poselb$ and $\poselb \posAleq \posela$ we get $\poselc \posAleq \posela$, by transitivity of $\posAleq$.
        Hence $\posela \simeq \poselc$.
    \end{enumerate}
    The second statement is a reformulation of the antisymmetry condition in \cref{def:poset}.
\end{proof}

Being an equivalence relation, $\simeq$ groups the elements of $\posAset$ into disjoint \emph{equivalence classes}.
The equivalence class of $\posela \setin \posAset$ is the set of all elements equivalent to it:
\begin{equation}
    [\posela] \definedas \makeset{ \poselb \setin \posAset \mid \posela \simeq \poselb }.
\end{equation}
Each element belongs to exactly one equivalence class, and two elements have the same class exactly when they are equivalent: $[\posela] = [\poselb]$ if and only if $\posela \simeq \poselb$.
In other words, the equivalence classes form a \SY{partition} of $\posAset$ (\cref{def:partition}).
The set of all equivalence classes is called the \emph{quotient set} and is denoted by $\posAset/_{\simeq}$.

We would like to order the equivalence classes by comparing their elements.
For this to make sense, the result must not depend on which elements of the classes we compare.
The following lemma shows that this is the case.

\begin{lemma}
    \label{lem:preorder-order-on-classes}
    Let $\posA$ be a \SY{preorder}, and let $\posela, \posela', \poselb, \poselb' \setin \posAset$ with $\posela \simeq \posela'$ and $\poselb \simeq \poselb'$.
    Then $\posela \posAleq \poselb$ if and only if $\posela' \posAleq \poselb'$.
\end{lemma}
\begin{proof}
    Suppose $\posela \posAleq \poselb$.
    Since $\posela' \posAleq \posela$ and $\poselb \posAleq \poselb'$, two applications of transitivity give $\posela' \posAleq \poselb'$.
    The converse follows in the same way, exchanging the roles of the primed and unprimed elements.
\end{proof}

\begin{ctdefinition}[Quotient poset of a preorder]
    \label{def:preorder-quotient}
    Let $\posA$ be a \SY{preorder}.
    The \maindef{quotient poset} of $\posA$ is the tuple $\posA/_{\simeq} = \tupp{\posAset/_{\simeq}, \posleq}$ given by:

    \constit
    \begin{enumerate}
        \item The carrier set $\posAset/_{\simeq}$ of equivalence classes;
        \item The order relation defined by
              \begin{equation}
                  [\posela] \posleq [\poselb]
                  \quad \text{if and only if} \quad
                  \posela \posAleq \poselb.
              \end{equation}
    \end{enumerate}
    By \cref{lem:preorder-order-on-classes}, this does not depend on the choice of the elements $\posela$ and $\poselb$ in their classes.
\end{ctdefinition}

\begin{lemma}
    \label{lem:preorder-quotient-is-poset}
    Let $\posA$ be a \SY{preorder}.
    Then $\posA/_{\simeq}$ is a \SY{poset}.
\end{lemma}
\begin{proof}
    Reflexivity and transitivity of $\posleq$ follow directly from the corresponding properties of $\posAleq$.
    For antisymmetry, suppose $[\posela] \posleq [\poselb]$ and $[\poselb] \posleq [\posela]$.
    By definition, this means $\posela \posAleq \poselb$ and $\poselb \posAleq \posela$, that is, $\posela \simeq \poselb$.
    Hence $[\posela] = [\poselb]$.
\end{proof}

Passing from $\posA$ to $\posA/_{\simeq}$ loses no information about the order: for all $\posela, \poselb \setin \posAset$, we have $\posela \posAleq \poselb$ if and only if $[\posela] \posleq [\poselb]$.
The only thing we forget is \emph{which} different elements were equivalent to each other.

Sometimes it is more convenient to stay inside the original set $\posAset$, rather than working with sets of elements.
Then we can instead choose one element from each equivalence class, which serves as a ``representative'' of its class.

\begin{ctdefinition}[Skeleton of a preorder]
    \label{def:preorder-skeleton}
    Let $\posA$ be a \SY{preorder}.
    A \maindef{skeleton} of $\posA$ is:

    \constit
    \begin{enumerate}
        \item A subset $\setS \setsubseteq \posAset$.
    \end{enumerate}

    \condit
    \begin{enumerate}
        \item $\setS$ contains exactly one element from each equivalence class of $\simeq$:
              for every $\posela \setin \posAset$ there is exactly one $\poselb \setin \setS$ with $\posela \simeq \poselb$.
    \end{enumerate}
    We equip $\setS$ with the restriction of $\posAleq$: for $\posela, \poselb \setin \setS$, we have $\posela \posleqof{\setS} \poselb$ if and only if $\posela \posAleq \poselb$.
\end{ctdefinition}

\begin{lemma}
    \label{lem:preorder-skeleton}
    Let $\posA$ be a \SY{preorder} and let $\setS$ be a \SY{skeleton} of $\posA$.
    Then $\tup{\setS, \posleqof{\setS}}$ is a \SY{poset}.
    Moreover, the function
    \begin{equation}
        \setS \sto \posAset/_{\simeq},
        \qquad
        \posela \mapsto [\posela]
    \end{equation}
    is a bijection, and for all $\posela, \poselb \setin \setS$ we have $\posela \posleqof{\setS} \poselb$ if and only if $[\posela] \posleq [\poselb]$.
\end{lemma}
\begin{proof}
    Reflexivity and transitivity of $\posleqof{\setS}$ are inherited from $\posAleq$.
    For antisymmetry, suppose $\posela, \poselb \setin \setS$ with $\posela \posleqof{\setS} \poselb$ and $\poselb \posleqof{\setS} \posela$.
    Then $\posela \simeq \poselb$, so $\posela$ and $\poselb$ are elements of $\setS$ from the same equivalence class.
    Since $\setS$ contains only one element from each class, $\posela = \poselb$.

    The function $\posela \mapsto [\posela]$ is surjective, because every equivalence class contains an element of $\setS$.
    It is injective, because $[\posela] = [\poselb]$ means $\posela \simeq \poselb$, which implies $\posela = \poselb$ as just shown.
    Finally, $\posela \posleqof{\setS} \poselb$ holds if and only if $\posela \posAleq \poselb$, which holds if and only if $[\posela] \posleq [\poselb]$ by \cref{def:preorder-quotient}.
\end{proof}

In general, a preorder has many different skeletons, one for each way of choosing the representatives.
However, \cref{lem:preorder-skeleton} tells us that each skeleton is a copy of the quotient poset $\posA/_{\simeq}$, in which each class $[\posela]$ is simply ``renamed'' by its chosen representative.
(In the language of \cref{def:order-isomorphism}, which we will meet later, every skeleton is order isomorphic to $\posA/_{\simeq}$.)
For this reason one often speaks of ``the'' skeleton of a preorder.
Note that if $\posA$ is already a \SY{poset}, then every equivalence class consists of a single element, and the only skeleton of $\posA$ is $\posAset$ itself.

\begin{example}[Comparing by price]
    \label{ex:skeleton-price}
    Suppose that a shop sells five chocolate bars $c_1, \dots, c_5$ with the following prices (in CHF):
    \begin{center}
        \begin{tabular}{c|ccccc}
            bar & $c_1$ & $c_2$ & $c_3$ & $c_4$ & $c_5$ \\
            \hline
            price & 2 & 3 & 3 & 5 & 2
        \end{tabular}
    \end{center}
    Let $\posAset = \makeset{c_1, c_2, c_3, c_4, c_5}$, and compare bars by price: $c_i \posAleq c_j$ if and only if the price of $c_i$ is less than or equal to the price of $c_j$.
    This is a \SY{preorder}, because the usual order on prices is reflexive and transitive.
    It is not a \SY{poset}, because $c_1 \posAleq c_5$ and $c_5 \posAleq c_1$, although $c_1 \neq c_5$.

    Two bars are equivalent exactly when they have the same price.
    Hence there are three equivalence classes,
    \begin{equation}
        [c_1] = [c_5] = \makeset{c_1, c_5},
        \qquad
        [c_2] = [c_3] = \makeset{c_2, c_3},
        \qquad
        [c_4] = \makeset{c_4},
    \end{equation}
    and the quotient poset $\posA/_{\simeq}$ is the \SY{totally ordered set} $[c_1] \posleq [c_2] \posleq [c_4]$.
    It is a copy of the set of prices that occur, $\makeset{2, 3, 5}$, with the usual order.
    The sets $\makeset{c_1, c_2, c_4}$ and $\makeset{c_5, c_3, c_4}$ are two of the four possible skeletons of $\posA$: in each, we keep one bar for every price.

    The same happens whenever a preorder is defined by comparing values of a function $\mapa \colon \setA \sto \reals$, via $\ela \posleq \elb$ if and only if $\mapa(\ela) \leq \mapa(\elb)$.
    Then $\ela \simeq \elb$ holds if and only if $\mapa(\ela) = \mapa(\elb)$, which is the equivalence relation $\equivalparam{\mapa}$ that we met when discussing equivalence relations, and the quotient poset is a copy of the image of $\mapa$ with the usual order of \reals.
\end{example}

\begin{marginfigure}
    \centering
    \subfloat[
        The directed graph of \cref{ex:skeleton-reachability}; the dashed boxes are the equivalence classes.
        \label{fig:skeleton-reachability-graph}
    ]{
        \includesag{40_reachability_skeleton_graph}
    }
    \\
    \subfloat[
        The quotient poset.
        \label{fig:skeleton-reachability-quotient}
    ]{
        \includesag{40_reachability_skeleton_quotient}
    }
    \caption{Passing from the reachability preorder of a directed graph to a poset.}
    \label{fig:skeleton-reachability}
\end{marginfigure}

\begin{example}[Reachability]
    \label{ex:skeleton-reachability}
    Consider the directed graph in \cref{fig:skeleton-reachability-graph}, with nodes $a, b, c, d, e$ and edges
    \begin{equation}
        a \to b, \quad
        b \to a, \quad
        b \to c, \quad
        c \to d, \quad
        d \to e, \quad
        e \to c,
    \end{equation}
    and let $\posA$ be the \SY{preorder} on the set of nodes given by reachability: $\ela \posAleq \elb$ if and only if there is a path from $\ela$ to $\elb$.
    As we saw in \cref{ex:reachabilityposet}, this is not a \SY{poset}, because the graph contains cycles.

    Two nodes are equivalent exactly when each can be reached from the other.
    Here, $a \simeq b$ and $c \simeq d \simeq e$, so there are two equivalence classes,
    \begin{equation}
        [a] = \makeset{a, b},
        \qquad
        [c] = \makeset{c, d, e}.
    \end{equation}
    Since there is a path from $a$ to $c$, but none from $c$ to $a$, the quotient poset $\posA/_{\simeq}$ has two elements, with $[a] \posleq [c]$, as shown in \cref{fig:skeleton-reachability-quotient}.
    A skeleton is obtained by choosing one node from each class, for example $\makeset{a, c}$ or $\makeset{b, e}$; there are $2 \cdot 3 = 6$ skeletons in total.

    In general, the equivalence classes of the reachability preorder of a directed graph are called the \emph{strongly connected components} of the graph.
    So, instead of restricting to acyclic graphs as in the solution of \cref{ex:reachabilityposet}, we can obtain a poset from any directed graph by collapsing each strongly connected component to a single element.
\end{example}
